%!TEX root=../../note
\subsection{Truth Table and Negation}
\begin{definition}
    The truth value for "not $A$", written symbolically as $\neg A$, is defined by the truth table given below. 

    \[
    \begin{array}{c|c}
    A & \neg A \\
    \hline
    \mathrm{T} & \mathrm{F} \\
    \mathrm{F} & \mathrm{T}
    \end{array}
    \]

    We also refer to $\neg A$ as the negation of $A$. 
\end{definition}

We can use the truth table to prove that an argument is equal to another argument. 

\begin{example}
    Given the truth table for $A, \neg A, \neg \paren{\neg A}$, 
    \[
        \begin{array}{c|c|c}
        A & \neg A & \neg(\neg A) \\
        \hline
        \mathrm{T} & \mathrm{F} & \mathrm{T} \\
        \mathrm{F} & \mathrm{T} & \mathrm{F}
        \end{array}
    \]

    This truth table serves as a proof that $A \equiv \neg \paren{\neg A}$. 
\end{example}

\begin{remark}
    "A" is considered a statement variable. 
\end{remark}

\subsection{Conjunction and Disjunction}
\begin{definition}
    [Compound statement]
    A compound statement is a statement composed of several individual statements, each of which is called a component statement. 
\end{definition}

\begin{definition}
    The truth value for $A \wedge B$, is defined by the truth table 

    \[
    \begin{array}{c|c|c}
    A & B & A \wedge B \\
    \hline
    \mathrm{T} & \mathrm{T} & \mathrm{T} \\
    \mathrm{T} & \mathrm{F} & \mathrm{F} \\
    \mathrm{F} & \mathrm{T} & \mathrm{F} \\
    \mathrm{F} & \mathrm{F} & \mathrm{F}
    \end{array}
    \]

    We also refer to $A \wedge B$ as the conjunction of $A$ and $B$. 
\end{definition}

If we have three component statement, we need $2^3$ rows to cover all the possibility. 

\begin{example}
    The compound statement 
    \begin{equation*}
        \paren{2 + 3 = 6} \wedge (\pi + 2 > 5)
    \end{equation*}

    is false, since the statement $2 + 3 = 6$ is false. 
\end{example}

\begin{definition}
    The truth value for $A \text{ or } B$, written as $A \vee B$, is defined by the truth table

    \[
    \begin{array}{c|c|c}
    A & B & A \vee B \\
    \hline
    \mathrm{T} & \mathrm{T} & \mathrm{T} \\
    \mathrm{T} & \mathrm{F} & \mathrm{T} \\
    \mathrm{F} & \mathrm{T} & \mathrm{T} \\
    \mathrm{F} & \mathrm{F} & \mathrm{F}
    \end{array}
    \]

    We also refer to $A \vee B$ as the \textbf{disjunction} of $A$ and $B$. 
\end{definition}

\begin{remark}
    The disjunction $A \vee B$ is the \emph{inclusive} or: it is true whenever at least one of $A$, $B$ is true, including the case where both are true. This differs from the \emph{exclusive} or (xor), written $A \oplus B$, which is true when exactly one of $A$, $B$ is true and false when both are true. Thus $A \vee B$ and $A \oplus B$ agree in every case except $A = B = \mathrm{T}$, where $A \vee B$ is true but $A \oplus B$ is false.
\end{remark}

\begin{example}
    The Compound statement: 
    \begin{equation*}
        (2 + 3 = 6) \vee (\pi + 2 \geq 5),
    \end{equation*}
    is true since one of the statements is true. 

\end{example}

\begin{example}
    Let us use a truth table to prove that 
    \begin{equation*}
        \neg (A \vee B) \equiv (\neg A \wedge \neg B)
    \end{equation*}

    \[
    \begin{array}{c|c|c|c|c|c|c}
    A & B & A\vee B & \neg(A\vee B) & \neg A & \neg B & \neg A\wedge\neg B\\
    \hline
    \mathrm{T} & \mathrm{T} & \mathrm{T} & \mathrm{F} & \mathrm{F} & \mathrm{F} & \mathrm{F}\\
    \mathrm{T} & \mathrm{F} & \mathrm{T} & \mathrm{F} & \mathrm{F} & \mathrm{T} & \mathrm{F}\\
    \mathrm{F} & \mathrm{T} & \mathrm{T} & \mathrm{F} & \mathrm{T} & \mathrm{F} & \mathrm{F}\\
    \mathrm{F} & \mathrm{F} & \mathrm{F} & \mathrm{T} & \mathrm{T} & \mathrm{T} & \mathrm{T}
    \end{array}
    \]
\end{example}

A similar equivalence is 
\begin{equation*}
    \neg (A \wedge B) \equiv \paren{\paren{\neg A} \vee \paren{\neg B}}. 
\end{equation*}

These things together are called DeMorgan's Law. 
\begin{theorem}
    [DeMorgan's Law]
    For statement variable $A$ and $B$, we have 
    \begin{align*}
         \neg (A \vee B) \equiv (\neg A \wedge \neg B) \\
        \neg (A \wedge B) \equiv (\neg A \vee \neg B)
    \end{align*}
\end{theorem}

\begin{proposition}
    For statements $A, B, C$, the following rules hold for the logical operations $\wedge$ and $\vee$: 

    \textbf{Commutative Laws}
    \begin{align*}
        A \vee B \equiv B \vee A \\
        B \wedge A \equiv A \wedge B
    \end{align*}

    \textbf{Associative Laws}
    \begin{align*}
        A \vee (B \vee C) \equiv (A \vee B) \vee C \\ 
        A \wedge (B \wedge C) \equiv (A \wedge B) \wedge C
    \end{align*}

    \textbf{Distributive Laws}
    \begin{align*}
        A \vee (B \wedge C) \equiv (A \vee B) \wedge (A \vee C) \\ 
        A \wedge (B \vee C) \equiv (A \wedge B) \vee (A \wedge C)
    \end{align*}

\end{proposition}

